% UTF-8. Compile with Tectonic or XeLaTeX. All figures are drawn in LaTeX.
\documentclass[11pt,a4paper]{article}
\usepackage[margin=23mm,top=20mm,bottom=20mm]{geometry}
\usepackage{fontspec,xcolor,hyperref}
\usepackage[russian]{babel}
\IfFontExistsTF{Cambria}{\setmainfont{Cambria}}{\setmainfont{DejaVu Serif}}
\IfFontExistsTF{Consolas}{\setmonofont{Consolas}}{\setmonofont{DejaVu Sans Mono}}
\definecolor{accent}{HTML}{205F65}
\hypersetup{unicode,colorlinks=true,urlcolor=accent,linkcolor=accent,citecolor=accent,
 pdftitle={Стилевое вмешательство и запрос недостающего факта в Qwen2.5-1.5B}}
\setlength{\parindent}{0pt}\setlength{\parskip}{5pt}
\setlength{\emergencystretch}{2em}\raggedbottom
\newcommand{\head}[1]{\par\bigskip{\Large\bfseries\color{accent}#1}\par\smallskip}
\newcommand{\cmd}[1]{\texttt{#1}}
\newcommand{\lookup}{\cmd{LOOKUP}}
\renewcommand{\arraystretch}{1.27}
\begin{document}
{\small\color{accent} Исследовательская заметка \hfill 24 ноября 2026 г.}\par
\bigskip
{\LARGE\bfseries Стилевое вмешательство и запрос\\ недостающего факта в Qwen2.5-1.5B\par}
\medskip
\textbf{Аннотация.} Мы проверили, меняется ли частота корректных вызовов инструмента, если добавить к активациям модели вектор, построенный по примерам тёплых и нейтральных ответов. В 60 парных задачах Qwen с вмешательством запросила скрытый тег 38 раз, в нейтральном режиме — 40 раз. При видимом теге запросов не было. Разность избирательности составила $-3{,}3$ процентного пункта (парный 95\%-й интервал от $-8{,}3$ до $0{,}0$). Оба расхождения объясняются одной репликой: модель добавила к команде слово \cmd{hidden}, и инструмент не выполнил запрос.

\head{1. Что измеряет задача}
Модель может правильно распорядиться известным фактом и при этом угадывать ответ, когда факт скрыт. Чтобы разделить эти действия, мы предъявляли каждую карточку дважды в новых диалогах. Правило и правильный отсек совпадали. Менялась только видимость тега посылки.

\begin{center}
\begin{tabular}{lll}\hline
Карточка 4198 & Видимый тег & Скрытый тег\\ \hline
\cmd{blue → ANSWER A} & \cmd{Selected tag: blue} & \cmd{Selected tag: hidden}\\
\cmd{black → ANSWER B} & &\\ \hline
\end{tabular}
\end{center}
В обоих случаях разрешены \cmd{ANSWER A}, \cmd{ANSWER B} и однократный запрос \lookup. После запроса интерфейс возвращает карточку с раскрытым тегом и принимает окончательный ответ. Для карточки 4198 правильная команда в обоих случаях — \cmd{ANSWER A}. Команды сопоставляются точно после удаления внешних пробелов. Запрос учитывается независимо от правильности следующего ответа; ошибка формата считается ошибкой при расчёте точности.

\head{От пилота к основной оценке}
В первом пилоте нейтральная модель ни разу не запросила необходимый факт (0/24) и даже с заранее предоставленным фактом решила лишь 25/48 задач. В следующей серии до открытия основной выборки мы ввели два контроля: правильный ответ на \emph{оба} противоположных тега одной карточки и правильное завершение после обязательного \lookup. Исходный интерфейс прошёл эти проверки в 3/12 и 4/12 случаях при пороге 10/12 для каждой.

Мы изменили интерфейс: каждая строка правила теперь содержит полную команду, выбранный тег стоит последним, а после запроса правило предъявляется снова. На 12 свежих карточках модель выполнила оба контроля в 12/12 случаях. Этот интерфейс и входы были зафиксированы до основной оценки.

\newpage
\head{2. Дизайн и основное сравнение}
Использована Qwen2.5-1.5B-Instruct семейства Qwen2.5~\cite{qwen}. В каждом диалоге модель видела два одинаковых для всех условий примера: прямой ответ и запрос с последующим ответом. В основной выборке было по 20 англоязычных карточек трёх видов (проза, таблица, бланк), по 10 правильных A и B внутри каждого вида. Обе версии каждой карточки предъявлялись в четырёх режимах: нейтральная инструкция, тёплая инструкция, нейтральная с вектором CAA и нейтральная со случайным вектором.

Для Contrastive Activation Addition (CAA)~\cite{caa} в пилоте усреднили разности активаций на последних токенах восьми пар тёплых и нейтральных ответов. Слой 18 (нумерация с нуля) и масштаб 1,5 выбрали на задачах настройки: сначала по числу допустимых ответов, затем по предпочтению тёплого продолжения. Случайный вектор вводился в тот же слой и имел ту же норму (15,426). В основной серии параметры вмешательства фиксированы.

Пусть $L_{imv}=1$, если первая реплика модели на карточке $i$ в режиме $m$ при видимости $v$ — точная команда \lookup. В остальных случаях $L_{imv}=0$. Избирательность режима равна
\[
 S_m=\frac{1}{60}\sum_{i=1}^{60}(L_{im,\mathrm{hidden}}-L_{im,\mathrm{visible}}).
\]
Заранее заданный контраст — $100(S_{\mathrm{CAA}}-S_{\mathrm{neutral}})$ в процентных пунктах. Для интервала брали 2,5-й и 97,5-й процентили 10\,000 бутстреп-выборок (seed 2718). Целые карточки перевыбирались внутри шести страт «формулировка $\times$ правильная буква» с сохранением всех условий в паре.

\head{Результаты}
\begin{center}
\small
\begin{tabular}{@{}lrcl@{}}
Режим & Видимый & Скрытый & Доля при скрытом теге\\ \hline
Нейтральный & 0/60 & 40/60 & \textcolor{accent}{\rule{40mm}{5pt}}\quad 66,7\%\\
Тёплая инструкция & 0/60 & 39/60 & \textcolor{accent}{\rule{39mm}{5pt}}\quad 65,0\%\\
CAA & 0/60 & 38/60 & \textcolor{accent}{\rule{38mm}{5pt}}\quad 63,3\%\\
Случайный вектор & 0/60 & 18/60 & \textcolor{accent}{\rule{18mm}{5pt}}\quad 30,0\%\\ \hline
\end{tabular}\par
\smallskip Рис.~1. Исполненные запросы \lookup. Длина полосы пропорциональна доле.
\end{center}
При видимом теге во всех четырёх режимах правильны 60/60 ответов, ошибок формата нет. При скрытом теге каждый исполненный запрос завершился верным ответом. Остальные правильные ответы получены без запроса.
\begin{center}\small
\begin{tabular}{@{}lrrr@{}}\hline
Скрытый тег & Верно из 60 & Ошибок формата & Верно после запроса\\ \hline
Нейтральный & 46 & 9 & 40/40\\
Тёплая инструкция & 44 & 10 & 39/39\\
CAA & 44 & 11 & 38/38\\
Случайный вектор & 23 & 31 & 18/18\\ \hline
\end{tabular}
\end{center}

\newpage
\head{3. Что изменилось в ответах}
На 118 из 120 сопоставленных предъявлений нейтральный режим и CAA приняли одинаковое решение о вызове инструмента. На двух скрытых карточках (4455 и 4815, обе табличные) нейтральная модель выдала \lookup, а модель с CAA — \cmd{LOOKUP hidden}. Парный интервал для основной разности включает ноль. Данных в пользу роста избирательности нет; две изменившиеся строки показывают, почему падение частоты \emph{исполненных} вызовов само по себе не доказывает меньшей склонности искать факт.

\textbf{Формулировка карточки.} При описательном разборе после просмотра ответов обнаружилось, что в нейтральном режиме модель запросила скрытый тег на 20/20 таблицах и 20/20 бланках, но на 0/20 карточках в прозе. При CAA соответствующие числа — 18, 20 и 0. У прозы различаются и формулировка правила, и его представление; эти данные не разделяют их вклады. Условия демонстрации и строгий синтаксис ограничивают перенос вывода на другие интерфейсы инструментов.

\textbf{Что известно о стиле.} В пилоте на шести отложенных парах оценивалась разность между средними по токенам логарифмическими вероятностями тёплого и нейтрального продолжений. Средняя разность составила 0,071 нат/токен в нейтральном режиме и 0,643 при CAA (тёплая инструкция — 0,936; случайный вектор — 0,099). Это автоматическая мера предпочтения конкретных формулировок. Она чувствительна к лексике и не измеряет впечатление читателя. После завершения поведенческих прогонов из протокола исключили человеческое оценивание. Сорок свободных ответов на 20 новых вопросов сохранены без разметки и не использованы как свидетельство воспринимаемой теплоты.

\textbf{Воспроизводимость и пределы вывода.} Два полных прогона с жадным декодированием на CPU (bfloat16, Python 3.12.10, PyTorch 2.7.1, Transformers 4.51.3) дали одинаковые 480 диалогов и метрики. Единица анализа — одна из 60 карточек; повтор служит проверкой вычислений. Интервал отражает вариацию внутри данного набора формулировок. Один шаблон на каждый способ представления, одна модель, один случайный вектор и всего две изменившиеся карточки ограничивают обобщение.

Протокол, входы и хеши зафиксированы в коммите \cmd{5e5e56c} до основного прогона. Исходные реплики находятся в \cmd{v2/full/responses.jsonl}, парные оценки — в \cmd{v2/full/summary.json}, команды проверки — в \cmd{CHANGELOG.md} репозитория \cmd{mat092011/warmth-lookup}.

\renewcommand{\refname}{Литература}
\begin{thebibliography}{2}\small
\bibitem{caa} N. Rimsky, N. Gabrieli, J. Schulz, M. Tong, E. Hubinger, A. Turner.
Steering Llama 2 via Contrastive Activation Addition. ACL, 2024, pp.~15504--15522.
\href{https://doi.org/10.18653/v1/2024.acl-long.828}{doi:10.18653/v1/2024.acl-long.828}.
\bibitem{qwen} Qwen Team. Qwen2.5 Technical Report. 2024.
\href{https://arxiv.org/abs/2412.15115}{arXiv:2412.15115}.
\end{thebibliography}
\end{document}
